%% Stage report: DeReFusion reproduction + operator-specialization diagnostics
%% Compile: pdflatex -interaction=nonstopmode stage_report_en.tex  (twice)
%% Companion file: stage_report_zh.tex (independent, xelatex + xeCJK)
\documentclass[preprint,12pt]{elsarticle}

\usepackage{amssymb}
\usepackage{amsmath}
\usepackage{booktabs}
\usepackage{graphicx}
\usepackage{float}
\usepackage[T1]{fontenc}
\usepackage{lmodern}
\usepackage{microtype}
\usepackage{hyperref}

\journal{Stage report --- internal}

\begin{document}

\begin{frontmatter}

\title{Reproduction of DeReFusion and Diagnostics of Operator
Specialization for Financial Time Series\tnoteref{t1}}
\tnotetext[t1]{Reproduction of Hsieh and Chen (2026), Applied Soft Computing 203, 116252,
extended with controlled diagnostics of regime- and asset-dependent operator suitability.}

\author[nwpu]{Dongxu Jiang\corref{cor1}}
\ead{jiangdongxu04@gmail.com}
\cortext[cor1]{Corresponding author. Email: jiangdongxu04@gmail.com.}
\affiliation[nwpu]{organization={},
            city={},
            country={}}

\begin{abstract}
This report consolidates the work completed up to 2026-09-12 on the reproduction of
DeReFusion and on a sequence of controlled diagnostics that ask a single question: does the
structural state of the input determine which operator is more useful? We first reproduce the
core comparison of the original paper (DLinear base branch, LSTM--Transformer residual branch,
parameter-free additive fusion, all wrapped in RevIN) on the S\&P~500 and two further assets.
We then stratify forecast errors by causal relative volatility and find that the marginal value
of the nonlinear branch is \emph{asset-dependent} rather than universal: it is significantly
positive for the S\&P~500 (and seed-robust), directionally positive but statistically
inconclusive for ETHUSD, and significantly \emph{reversed} for BTCUSD. A capacity-sensitivity
experiment (hidden width $23\!\to\!128$, everything else frozen) shows that the earlier
``linear dominates everywhere'' reading was partly an artefact of a narrow nonlinear bottleneck:
with adequate capacity ETHUSD prefers the nonlinear operator in $36/36$ observations, while
BTCUSD never does. A capacity-matched proxy routing experiment (12 pre-registered structural
states, linear vs.\ MLP, identical protocol) finds no robust state-level operator specialization.
Finally, an exploratory sweep over nine assets shows interaction effects that span zero
(six negative, three positive; five statistically supported) and two leave-one-out-stable
associations (realized volatility and $|\mathrm{ACF}_1|$), which we report as a
\emph{candidate} asset-level structural regularity only. We state explicitly what the evidence
excludes (sample-level routing), what it supports (asset-level operator heterogeneity), and what
remains unresolved.
\end{abstract}

\begin{keyword}
reproduction \sep time series forecasting \sep operator specialization \sep
structural states \sep adaptive fusion \sep controlled comparison
\end{keyword}

\end{frontmatter}

\section{Scope and research questions}
\label{sec:intro}

The work reported here follows an evidence-first protocol: reproduce the reference result,
then ask one question at a time, with every grouping rule, threshold and operator definition
fixed \emph{before} looking at outcomes. Three questions structure the report.

\begin{itemize}
\item \textbf{Q-A (reproduction).} Does the published ordering of fusion strategies
(additive fusion $>$ linear base $>$ learnable gate) reproduce on independently fetched data?
\item \textbf{Q-B (regime).} Is the marginal contribution of the nonlinear branch conditioned on
the volatility regime of the input window?
\item \textbf{Q-C (specialization).} Does the \emph{structural state} of the input determine
which operator is more useful --- i.e.\ is there an experimental basis for routing?
\end{itemize}

All experiments run on CPU, use a frozen 70/10/20 chronological split, and report paired
bootstrap confidence intervals (4000 resamples) with paired win rates. No hyper-parameter was
tuned after observing test outcomes.

\section{Reproduction of the reference comparison}
\label{sec:repro}

The dataset of the original repository is not redistributable, so ten daily OHLC series
(2016--2025) were re-fetched from Yahoo Finance; row counts match the paper (2513 for indices
and equities, 2602 for FX, 3653 for BTCUSD, 2975 for ETHUSD). Table~\ref{tab:repro} reports the
S\&P~500 comparison at horizons $T\in\{1,24\}$ (seed 2021).

\begin{table}[!t]
\centering
\small
\caption{Reproduction of the core comparison (GSPC, seed 2021, normalized scale).}
\label{tab:repro}
\begin{tabular}{@{}lcccc@{}}
\toprule
Model & $T{=}1$ MSE & $T{=}24$ MSE & $T{=}24$ $R^2$ & Params \\
\midrule
DeReFusion (additive)      & 0.01647 & \textbf{0.06230} & \textbf{0.8691} & 28{,}436 \\
revin-DLinear (linear)     & \textbf{0.01577} & 0.07007 & 0.8528 & 4{,}664 \\
gatev1 (volatility-aware)  & ---     & 0.06767 & 0.8578 & 28{,}580 \\
gatev2 (learnable)         & 0.02526 & 0.07191 & 0.8489 & $\sim$28k \\
\bottomrule
\end{tabular}
\end{table}

At $T=24$ the published ordering reproduces: additive fusion outperforms the linear base branch
by $11.1\%$ in MSE, while both gate variants are worse than parameter-free addition. At $T=1$
the linear branch is marginally better, in line with the ``decomposition-then-simple-model''
behaviour of long-horizon linear forecasters. Extending the same protocol to two further assets
at $T=24$ gives DeReFusion $0.21797$ vs.\ DLinear $0.22322$ for BTCUSD ($+2.4\%$) and
$0.14602$ vs.\ $0.14876$ for ETHUSD ($+1.8\%$); the additive advantage is therefore ordered
GSPC $\gg$ BTCUSD $\approx$ ETHUSD.

\section{Volatility-regime stratification}
\label{sec:regime}

Each test window is assigned a regime from a \emph{causal} volatility measure computed on the
input window only: the standard deviation of log returns, and a de-trended version,
$RV^{rel}_t = RV_t / \mathrm{median}(RV_{t-120:t-1})$. Table~\ref{tab:regime} reports the
relative improvement of the nonlinear branch over the linear base per stratum (relative
volatility, the primary measure) plus the interaction effect
$\Delta_{high50} - \Delta_{low50}$.

\begin{table}[!t]
\centering
\small
\caption{Stratified performance, relative volatility ($\Delta = $ MSE$_{nonlin}-$MSE$_{lin}$;
negative favouring the nonlinear branch). $^\dagger$CI excludes zero.}
\label{tab:regime}
\begin{tabular}{@{}lcccc@{}}
\toprule
Asset & low 20\% & high 50\% & high 20\% & Interaction [95\% CI] \\
\midrule
GSPC (2021) & $+6.0\%$      & $+14.2\%^\dagger$ & $+21.6\%^\dagger$ & $-0.01205\,[-0.01545,-0.00876]^\dagger$ \\
GSPC (2022) & $-2.7\%$      & $+23.2\%^\dagger$ & $+29.6\%^\dagger$ & $-0.02248\,[-0.02719,-0.01792]^\dagger$ \\
ETHUSD      & $+16.7\%^\dagger$ & $+6.3\%^\dagger$ & $+14.2\%^\dagger$ & $-0.00943\,[-0.01881,+0.00033]$ \\
BTCUSD      & $+3.5\%$      & $-1.0\%$          & $-3.5\%$          & $+0.01502\,[+0.00615,+0.02378]^\dagger$ \\
\bottomrule
\end{tabular}
\end{table}

Two findings follow. First, the S\&P~500 effect is \emph{seed-robust}: both seeds give a
significantly negative interaction of the same order, and the calm-regime strata are
indistinguishable between branches. Second, BTCUSD reverses the sign significantly, so
``nonlinear advantage grows with volatility'' cannot be stated as a universal criterion.
A methodological check explains part of the earlier confusion: the raw volatility measure is
strongly confounded with time (correlation of $RV$ with the sample index: $+0.44$ for GSPC,
$-0.67$ for BTCUSD). The de-trended $RV^{rel}$ removes that trend; under it the spurious
``linear wins in the middle quintile'' of GSPC disappears. Reported figures therefore use
$RV^{rel}$ as the primary measure.

\section{Adaptive fusion (gating) does not pay}
\label{sec:gate}

The natural use of the regime structure would be a gate $\alpha(X)$ mixing the two branches.
The framework already provides exactly this variant, and we ran it on the asset where the
structure is strongest and seed-robust (GSPC, $T=24$): the volatility-aware gate obtains MSE
$0.06767$, i.e.\ $8.6\%$ \emph{worse} than the parameter-free additive fusion ($0.06230$) and
only $3.4\%$ better than the plain linear base ($0.07007$). Since the nonlinear branch is never
worse than the linear one in the calm strata, the optimal gate weight is essentially constant
at one, leaving no switching headroom for the gate to exploit. Together with the previously
reported negative results for the learnable gate, this motivates stopping the adaptive-fusion
line of work rather than building a router on top of it.

\section{Capacity sensitivity: the rigidity was partly an artefact}
\label{sec:capacity}

The proxy experiment of the next section compares a linear operator with a capacity-matched MLP
(9{,}240 vs.\ 9{,}431 parameters, $+2.1\%$). That matching forces a 384-dimensional input
through a 23-dimensional hidden layer --- a severe bottleneck. We therefore re-ran the
structural-state evaluation with the hidden width as the \emph{only} variable
($23 \to 64 \to 128$; data, features, splits, optimizer, learning rate, batch size, epochs,
early stopping, seeds and evaluation protocol unchanged). Table~\ref{tab:cap} reports, per
width and asset, the mean $\Delta$MSE over 12 structural states and 3 seeds.

\begin{table}[!t]
\centering
\small
\caption{Capacity sensitivity. Mean $\Delta$MSE (negative $=$ nonlinear better); in
parentheses the number of the 36 (state, seed) observations preferring the nonlinear operator.}
\label{tab:cap}
\begin{tabular}{@{}lcccc@{}}
\toprule
Hidden & MLP params & GSPC & BTCUSD & ETHUSD \\
\midrule
23  & 9{,}431  & $+0.0598$ (0/36)  & $+0.1783$ (0/36)  & $+0.0145$ (12/36) \\
64  & 26{,}200 & $+0.0250$ (7/36)  & $+0.1888$ (0/36)  & $\mathbf{-0.0281}$ (36/36) \\
128 & 52{,}376 & $+0.0061$ (8/36)  & $+0.2998$ (0/36)  & $\mathbf{-0.0321}$ (36/36) \\
\bottomrule
\end{tabular}
\end{table}

The three pre-registered checks give a consistent reading. (i) \emph{Sign reversals}: 25
state-level reversals between widths, of which 23 belong to ETHUSD (all three seeds) and 8 to
GSPC (all of them at seed 2023 only, i.e.\ seed-unstable); BTCUSD has none. (ii)
\emph{Spread across states}: the range of $\Delta$MSE across states \emph{shrinks} with capacity
for GSPC ($0.106 \to 0.060$) and ETHUSD ($0.089 \to 0.040$); the BTCUSD increase is driven by
the magnitude of a single high-error state. (iii) \emph{Per-seed direction consistency} falls
from $31/36$ to $29/36$ to $28/36$. Taken together: the ETHUSD preference for the nonlinear
operator is \emph{asset-wide} once capacity is sufficient --- all twelve states reverse, not a
subset --- so it is asset-level heterogeneity rather than state-level specialization. Accordingly
we record the earlier ``linear dominates everywhere'' statement as partly capacity-driven, and
we do not treat the capacity sweep as evidence of routing.

\section{Structure-aware routing: a minimal controlled proxy experiment}
\label{sec:routing}

\subsection{Design}
Twelve structural states were defined \emph{before} any outcome was inspected, using
causal window features (ACF at lags 1/5/10, a trend $t$-statistic, sign persistence, a jump
ratio with a fixed threshold $c=3\sigma$, the de-trended relative volatility, skewness and
kurtosis). Binary states use train-set medians as thresholds; a composite \emph{structure
strength} state and a $K{=}3$ $k$-means cross-check are also reported. The operators are a
linear map ($9{,}240$ parameters) and an MLP with one 23-unit hidden layer ($9{,}431$
parameters), trained under an identical protocol (same seed, optimizer, steps, batch,
learning rate, early stopping, preprocessing).

\subsection{Result}
For GSPC and BTCUSD the linear operator is better in \emph{every} one of the twelve states at
matched capacity (GSPC $\Delta$MSE between $+0.047$ and $+0.081$, significant in all three
seeds; BTCUSD between $+0.129$ and $+0.448$). Only ETHUSD shows between-state variation of both
signs (structure-weak states favouring the nonlinear operator by up to $-0.013$, structure-strong
states favouring the linear one by $+0.033$). No state shows a preference that survives the
capacity sweep. \textbf{Main finding: no robust operator specialization is established at
capacity-matched scale.} The one cross-asset invariant is that the two operators are
indistinguishable in calm regimes.

\section{Asset-level operator heterogeneity across nine assets}
\label{sec:assets}

Finally, the same frozen protocol ($T=24$, seed 2021) was run on seven additional assets to map
the \emph{sign} of the interaction effect. GSPC aggregates its two seeds into a single asset, so
$N=9$. Table~\ref{tab:assets} reports the interaction effect with its confidence interval, and
Table~\ref{tab:assoc} the exploratory association between asset-level structure features and the
interaction effect, including a leave-one-asset-out (LOO) check.

\begin{table}[!t]
\centering
\small
\caption{Asset-level interaction effects (relative volatility). $^\dagger$CI excludes zero.}
\label{tab:assets}
\begin{tabular}{@{}lccc@{}}
\toprule
Asset & $\Delta_{interaction}$ & 95\% CI & Supported? \\
\midrule
EURUSD & $-0.1371$ & $[-0.164,-0.111]$ & yes$^\dagger$ \\
GSPC   & $-0.0173$ & $[-0.027,-0.009]$ & yes$^\dagger$ \\
USDJPY & $-0.0122$ & $[-0.026,+0.002]$ & no \\
ETHUSD & $-0.0094$ & $[-0.019,+0.000]$ & no \\
DJI    & $-0.0084$ & $[-0.019,+0.003]$ & no \\
TM     & $-0.0081$ & $[-0.028,+0.011]$ & no \\
BTCUSD & $+0.0150$ & $[+0.006,+0.024]$ & yes$^\dagger$ \\
BABA   & $+0.0405$ & $[+0.022,+0.060]$ & yes$^\dagger$ \\
NVO    & $+0.0976$ & $[+0.046,+0.149]$ & yes$^\dagger$ \\
\bottomrule
\end{tabular}
\end{table}

\begin{table}[!t]
\centering
\small
\caption{Exploratory Spearman association of asset structure features with
$\Delta_{interaction}$ ($N=9$) and LOO stability.}
\label{tab:assoc}
\begin{tabular}{@{}lcccc@{}}
\toprule
Feature & $\rho$ & $p$ & LOO range & Stable? \\
\midrule
$|\mathrm{ACF}_1|$      & $+0.683$ & $0.042$ & $[+0.548,+0.810]$ & yes \\
Realized volatility     & $+0.633$ & $0.067$ & $[+0.476,+0.857]$ & yes \\
Kurtosis                & $+0.533$ & $0.139$ & $[+0.333,+0.643]$ & yes \\
Trend persistence       & $+0.100$ & $0.798$ & $[-0.286,+0.381]$ & no (sign flips) \\
Jump ratio              & ---      & ---     & --- & constant (0.0105) \\
\bottomrule
\end{tabular}
\end{table}

The sign of the effect spans zero: six assets are negative, three positive, with five
statistically supported. Two features are LOO-stable: assets with higher realized volatility and
stronger first-order autocorrelation show a \emph{less} favourable (or reversed) nonlinear
advantage in turbulent regimes. \textbf{These associations are exploratory and are reported as a
candidate regularity only} --- with nine assets, a single $p$-value is not decisive, and the
jump-ratio feature turns out to be constant across all assets at the median (one exceedance in
95 returns), i.e.\ it carries no discriminative information at the pre-registered threshold.
Trend persistence flips sign under LOO and is marked \emph{single-asset sensitive}.

\section{Diagnosis and decisions}
\label{sec:diagnosis}

\begin{table}[!t]
\centering
\small
\caption{Consolidated answers to the five diagnostic questions.}
\label{tab:q}
\begin{tabular}{@{}p{0.44\textwidth}p{0.46\textwidth}@{}}
\toprule
Question & Answer \\
\midrule
Q1 Robust sample-level routing evidence? &
No. Neither the routing proxy nor the capacity sweep produces a stable
state $\rightarrow$ operator mapping. \\
Q2 Genuine asset-level operator heterogeneity? &
Yes. At adequate capacity ETHUSD prefers the nonlinear operator in 36/36
observations, BTCUSD prefers linear at every capacity, GSPC is near parity. \\
Q3 Explainable by the existing structural features? &
Only as a candidate: realized volatility and $|\mathrm{ACF}_1|$ associate with
the interaction sign and survive LOO; the sign spans zero across nine assets. \\
Q4 Principal remaining uncertainty? &
Asset dependence and operator form --- not the structural representation, and
not routing, which is now largely excluded. \\
Q5 May a stronger nonlinear operator be reconsidered? &
Only as one candidate, tested against a capacity-matched MLP, and not by
extrapolating ``ETHUSD prefers nonlinear'' to any specific operator family. \\
\bottomrule
\end{tabular}
\end{table}

The practical decisions recorded here are: (i) stop the adaptive-fusion line --- a volatility
gate measures worse than parameter-free addition, and the structure it would exploit leaves no
switching headroom; (ii) do not build a neural router --- the state $\rightarrow$ operator
mapping is not established; (iii) treat operator choice as an asset-level configuration
question, which is what the evidence supports; and (iv) keep any structured nonlinear operator
conditional on a reproducible context and on a capacity-matched comparison.

\section{Limitations}
\label{sec:limitations}

\textbf{Operator fidelity.} The proxy operators (a linear map and an MLP over flattened
windows) differ from the branches of the framework, which combine a DLinear decomposition with
an LSTM--Transformer residual branch under RevIN. Consequently, the proxy results must not be
read as ``nonlinearity is useless''; within the framework the residual branch still outperforms
the linear base on the S\&P~500 at $T=24$ across both seeds, and those two statements are
consistent.
\textbf{Sample size.} The asset-level association rests on nine assets and one seed per asset
(two for GSPC); confidence intervals for association coefficients are wide and the analysis is
exploratory.
\textbf{Regime definition.} Regimes are volatility quantiles of a causal window measure; the
calm-regime strata are statistically indistinguishable between operators, so ``calm'' is an
acceptance result rather than a separate effect.
\textbf{Environment.} Runs are CPU-only; wall-clock times differ from published GPU numbers and
only relative comparisons are interpreted.

\section{Conclusion}
\label{sec:conclusion}

The reproduction reproduces the published ordering at $T=24$ and reproduces it robustly across
seeds. The diagnostics that follow refine, but do not overturn, that result: the marginal value
of the nonlinear branch is asset-dependent rather than regime-universal, the volatility gate
does not pay, and matched-capacity comparisons show that the earlier uniform preference for the
linear operator was partly a capacity artefact. No robust experimental basis for sample-level
routing was found, whereas genuine asset-level operator heterogeneity was. The next question the
evidence supports is therefore not ``how should a router choose?'' but ``in which reproducible
contexts is a stronger nonlinear operator worth its capacity?'', to be answered against a
capacity-matched baseline.

\appendix
\section{Reproducibility notes}
\label{app:repro}

All reproduction assets live under \texttt{reproduction/} in the project repository:
analysis scripts in \texttt{reproduction/analysis/}, batch runners in
\texttt{reproduction/batches/}, and per-setting results in \texttt{reproduction/results/}.
Representative commands:

\begin{verbatim}
# single run (CPU)
.venv\Scripts\python.exe run.py --task_name long_term_forecast --is_training 1 \
  --model_id GSPC_96_24 --model DeReFusion --data custom --root_path ./dataset/ \
  --data_path GSPC-2016-2025.csv --features MS --target Close --freq b \
  --seq_len 96 --label_len 48 --pred_len 24 --enc_in 4 --dec_in 4 --c_out 1 \
  --d_model 32 --moving_avg 25 --train_epochs 30 --batch_size 32 \
  --learning_rate 0.0001 --patience 5 --lradj cosine --rand_seed 2021 --no_use_gpu

# regime stratification (causal, primary measure = relative volatility)
.venv\Scripts\python.exe reproduction/analysis/analyze_volatility_regimes.py \
  --csv dataset/GSPC-2016-2025.csv --tag GSPC --seed 2021 --rv-mode relative
\end{verbatim}

Four environment facts cost real time and are recorded for reuse: the data pipeline imports
\texttt{patool}, \texttt{huggingface\_hub}, \texttt{sktime} and \texttt{datasets} unconditionally;
\texttt{joblib} must be pinned to \texttt{1.5.3}; CPU runs require \texttt{--no\_use\_gpu}; and
killing DataLoader worker processes (\texttt{spawn\_main}) deadlocks the parent trainer, so a
stalled run must be detected by process CPU time rather than by log output.

\begin{thebibliography}{9}
\bibitem{hsieh2026} Hsieh, C.-Y., Chen, Y.-C. (2026). DeReFusion: a decomposition-reconstruction
fusion framework for financial time series forecasting. \emph{Applied Soft Computing} 203, 116252.
\bibitem{zeng2023} Zeng, A., Chen, M., Zhang, L., Xu, Q. (2023). Are transformers effective for
time series forecasting? \emph{Proceedings of AAAI} 37(9), 11121--11128.
\end{thebibliography}

\end{document}
